\begin{lemma}[Parameters for host refinement]
For an integer $k\geq1$, an integer extraction cutoff $N$, and real
numbers $0<\eta<1$, $0<\theta$, there is $e>0$, depending only on
$k,\eta,\theta$, such that
\[
 e<\eta,\quad e<1,\quad e<1/(8k^2),\quad
 (4k^2+2)e\leq\theta,\quad 2ke\leq\theta.
\]
For every integer $r\geq0$ there is $n_0$ such that for all $n\geq n_0$,
$n\geq2k$, $\lfloor n/k\rfloor\geq N$, and
\[
 2rk(r+1)n\exp\left(-\frac{(e/(2k))^2}{2(1+e/(6k))}n\right)<1.
\]
\end{lemma}
\begin{proof}
Take $e=\min\{\eta/2,\theta/(8k^2),1/(16k^2)\}$.
All three entries are positive. Since $k\geq1$,
$4k^2+2\leq8k^2$ and $2k\leq8k^2$; these inequalities give
the two non-strict margins. The first entry gives $e<\eta<1$;
the last gives $e<1/(8k^2)$.
Set $c=(e/(2k))^2/(2(1+e/(6k)))>0$. For each fixed $r$, the
function $2rk(r+1)n\exp(-cn)$ tends to zero as $n$ tends to infinity.
For instance $x\exp(-x)\to0$ with $x=cn$, and the prefactor
becomes the fixed constant $2rk(r+1)/c$. Choose an integer $M$
above which this function is less than one. Then
$n_0=\max\{2k,kN,M\}$ gives all three assertions; $kN\leq n$
implies $N\leq\lfloor n/k\rfloor$ because $k>0$.
When $r=0$ the tail expression is zero for every $n$ and this
same construction still works. Thus $e$ is selected before $r$;
only the order cutoff can depend on $r$. The positive rate here is
the weaker hypergeometric rate actually used in host sampling.
\end{proof}
